%! Change in Arithmetic and Geometric Sequences — Unit 4, Lesson 4.1 — Lesson Plan
\documentclass[10pt]{article}
\usepackage{precalculus-article}
\usepackage{precalculus-boxes}
\usepackage{graphicx}
\graphicspath{{images/}}
\newcommand{\TallMath}[1]{$\displaystyle #1\rule[-1.4em]{0pt}{3.2em}$}
\newcommand{\CourseName}{Precalculus}
\newcommand{\MeetingLength}{55 minutes}
\newcommand{\UnitNumberName}{Unit 4: Exponential Functions \quad}
\newcommand{\LessonNumberName}{Lesson 4.1: Change in Arithmetic and Geometric Sequences}

\begin{document}
\begin{center}
    {\Large\bfseries \CourseName \\
    {\small\bfseries \UnitNumberName \LessonNumberName}}
\end{center}

\begin{tcolorbox}[colback=sky, colframe=navy, boxrule=0.9pt, arc=2mm,
    left=2mm, right=2mm, top=1.5mm, bottom=1.5mm]
    \textbf{Primary Objective:} Students will express arithmetic and geometric sequences as
    functions of the whole numbers using explicit formulas, identify common differences and
    common ratios, and compare how the two sequence types change across successive terms.
    \hfill\textit{\small Skills: AP 1.B, AP 3.A}
\end{tcolorbox}

\renewcommand{\arraystretch}{1.6}
\begin{skillbox}[Priority Ideas \& Skills]{goldbox}
\begin{tabularx}{\linewidth}{@{}X X@{}}
\textbf{AP Skills} &
\textbf{Key Understandings} \\
\midrule
\textbf{AP 1.B} --- Identify an appropriate mathematical rule or procedure based on the
classification of a given expression. &
\begin{itemize}[leftmargin=1.2em,itemsep=2pt,topsep=0pt]
  \item A sequence is a function from the whole numbers to the reals; its graph is discrete
        points, not a curve. (EK 2.1.A.1)
  \item Arithmetic sequences have a common difference $d$; general term $a_n = a_0 + dn$.
        (EK 2.1.A.2, 2.1.A.3)
\end{itemize} \\
\textbf{AP 3.A} --- Describe the characteristics of a function. &
\begin{itemize}[leftmargin=1.2em,itemsep=2pt,topsep=0pt]
  \item Geometric sequences have a common ratio $r$; general term $g_n = g_0 r^n$.
        (EK 2.1.B.1, 2.1.B.2)
  \item Arithmetic sequences increase equally per step; geometric sequences increase by a
        larger amount each successive step. (EK 2.1.B.3)
\end{itemize} \\
\end{tabularx}
\end{skillbox}

\begin{skillbox}[{Vocabulary, Concepts, \& Theorems}]{greenbox}
\renewcommand{\arraystretch}{1.4}
\begin{tabularx}{\linewidth}{@{} >{\leavevmode\bfseries\color{navy}}p{0.27\linewidth} X @{}}
Sequence            & A function whose domain is the whole numbers $\{0, 1, 2, \ldots\}$; its graph consists of discrete points. \\
Arithmetic sequence & A sequence in which consecutive terms differ by a constant value $d$ (the \textit{common difference}). \\
Geometric sequence  & A sequence in which consecutive terms have a constant ratio $r$ (the \textit{common ratio}). \\
Common difference   & The constant $d$ added to each term to produce the next term in an arithmetic sequence. \\
Common ratio        & The constant $r$ by which each term is multiplied to produce the next term in a geometric sequence. \\
General term / explicit formula & A formula giving the $n$th term directly: $a_n = a_0 + dn$ (arithmetic) or $g_n = g_0 r^n$ (geometric). \\
Recursive formula   & A formula defining each term in terms of the previous term: $a_n = a_{n-1} + d$ or $g_n = g_{n-1} \cdot r$. \\
\end{tabularx}
\end{skillbox}

\begin{skillbox}[Activate Prior Knowledge \& Spiral Review]{sky}
    Spiral review (warm-up) covers: evaluating a linear function $f(x) = 3x - 1$ at integer
    inputs and identifying the constant rate of change (Unit 1 linear rate of change), stating
    the domain of $f(x) = \sqrt{x}$ (function notation and domain), and evaluating
    $f(n) = 2 \cdot 3^n$ at $n = 0, 1, 2$ (exponential function notation as preview).
    These connect to sequences as functions with restricted whole-number domains.
\end{skillbox}

\begin{skillbox}[Hook]{sky}
Write the two sequences on the board:
\[
  2,\ 5,\ 8,\ 11,\ \ldots \qquad\qquad 2,\ 6,\ 18,\ 54,\ \ldots
\]
Ask: \textbf{``What is the pattern in each? How are they different?''}
After student responses: \textbf{``How many terms until each sequence exceeds 1000?''}
Have students estimate by extending the table. (Arithmetic: about 333 terms;
Geometric: about 6--7 terms.) This creates immediate motivation for explicit formulas.
\end{skillbox}

\begin{skillbox}[Lesson]{sky}
\begin{multicols}{2}
\textbf{Part 1 (15 min): Arithmetic Sequences.}

Define: a sequence is a function from whole numbers to reals. The graph of
$2, 5, 8, 11, \ldots$ is a set of \textit{discrete points} (not a line).

Common difference: $d = a_n - a_{n-1}$ is constant. For $2, 5, 8, 11, \ldots$, $d = 3$.

Explicit formula: $a_n = a_0 + dn$. Starting at $a_0 = 2$:
$a_n = 2 + 3n$.
Check: $a_0 = 2$, $a_1 = 5$, $a_3 = 11$. \checkmark

General form when $a_k$ is known: $a_n = a_k + d(n - k)$.

Pose: ``How many terms until the sequence first exceeds 1000?''
Solve $2 + 3n > 1000 \Rightarrow n > 332.6$, so $n = 333$.

\columnbreak

\textbf{Part 2 (15 min): Geometric Sequences.}

Common ratio: $r = g_n / g_{n-1}$ is constant. For $2, 6, 18, 54, \ldots$, $r = 3$.

Explicit formula: $g_n = g_0 r^n$. Starting at $g_0 = 2$:
$g_n = 2 \cdot 3^n$.
Check: $g_0 = 2$, $g_1 = 6$, $g_3 = 54$. \checkmark

General form: $g_n = g_k \cdot r^{(n-k)}$.

Solve: ``When does $2 \cdot 3^n > 1000$?'' --- $3^n > 500$; since $3^5 = 243$
and $3^6 = 729 > 500$, $n = 6$. Only 6 steps vs.\ 333!

\textbf{Part 3 (10 min): Comparison.}

Arithmetic adds a constant each step; geometric multiplies by a constant.
On a table: show both sequences at $n = 0, 1, 2, \ldots, 6$.
Key insight (EK 2.1.B.3): increasing geometric sequences eventually grow
faster than any arithmetic sequence --- the gap widens each step.
\end{multicols}
\end{skillbox}

\begin{skillbox}[Active Monitoring]{redbox}
\begin{itemize}[leftmargin=1.2em,itemsep=2pt,topsep=0pt]
  \item After defining arithmetic sequences: cold-call ``Is a sequence a function? What is the domain?''
        Expect: domain is whole numbers $\{0, 1, 2, \ldots\}$, codomain is reals.
  \item During Part 1: watch for students who write $a_n = a_1 + d(n-1)$ (1-indexed) rather than
        $a_n = a_0 + dn$ (0-indexed). Clarify both conventions; this course uses index starting at 0.
  \item After Part 2: cold-call ``If $g_0 = 5$ and $r = -2$, does the sequence always increase?''
        (No; alternates sign.) Preview: geometric with $r < 0$ or $0 < r < 1$.
  \item During Part 3: ensure students articulate \textit{why} geometric sequences outpace arithmetic
        ones (multiplicative vs.\ additive change), not just that they do.
\end{itemize}
\end{skillbox}

\begin{skillbox}[Group Work \& Differentiation]{redbox}
\begin{itemize}[leftmargin=1.2em,itemsep=2pt,topsep=0pt]
  \item \textbf{Tier R (Remediate):} Given a sequence explicitly, identify type (arithmetic/geometric),
        find $d$ or $r$, write the general formula, evaluate the 10th term.
  \item \textbf{Tier A (Apply):} Given two terms (not the first), find $d$ or $r$ algebraically and
        write the explicit formula. Solve a contextual question about when a sequence exceeds a value.
  \item \textbf{Tier E (Extend):} Work from the formula $g_n = g_k r^{(n-k)}$ with a known interior
        term to find $g_0$; determine when a population first exceeds a target using the formula.
\end{itemize}
\end{skillbox}

\begin{skillbox}[Individual Work \& Assessment]{redbox}
Exit ticket (3 questions, 1 page): (1) write the explicit formula $a_n$ for $7, 11, 15, 19, \ldots$;
(2) write the explicit formula $g_n$ for $3, 6, 12, 24, \ldots$;
(3) true or false: both sequences in \#1 and \#2 are functions from the whole numbers to the reals.

Sort by performance: students missing \#1/\#2 need formula practice (LO 2.1.A or 2.1.B);
students missing \#3 need reinforcement that sequences \textit{are} functions (EK 2.1.A.1).
\end{skillbox}

\begin{skillbox}[Reinforcement \& Extension]{goldbox}
Homework (6 problems): identify arithmetic vs.\ geometric from a table and find $d$ or $r$;
write the general formula from two given terms; contextual cell-phone plan (arithmetic);
find which term first exceeds a value (arithmetic case); AP-style MC on $a_n = 4 + 3n$;
extension: decide whether a given sequence is arithmetic, geometric, or neither and explain.

Preview of Lesson 2.2: exponential functions --- connecting geometric sequences ($g_n = g_0 r^n$)
to the continuous function $f(x) = a \cdot b^x$ and extending the domain from whole numbers to all reals.
\end{skillbox}

\end{document}
